\chapter{Missão na era do medo}
\noindent\textit{Vinheta.} Uma igreja pequena decide ``ir para a internet''. Sem orçamento, começam com um celular, um tripé barato e coragem. Em seis meses, não viralizaram; mas três famílias começaram um grupo no bairro, um jovem foi batizado e uma senhora voltou à fé. \textbf{Evangelismo digital} não é palco: é \textbf{ponte}.

\section{Mapa bíblico}
\begin{itemize}[leftmargin=*,noitemsep]
  \item \textbf{A Palavra corre} (2~Ts~3:1): pedir que o evangelho tenha livre curso.
  \item \textbf{Ser astutos e simples} (Mt~10:16): prudência com pureza.
  \item \textbf{Testemunhas até os confins} (At~1:8): poder do Espírito para atravessar fronteiras culturais e tecnológicas.
  \item \textbf{Vida visível} (1~Pe~3:15): responder com mansidão e respeito, com boa consciência.
\end{itemize}

\begin{nicbox}
\textbf{NIC aplicado}\\
\textbf{Narrativa}: contamos a \textbf{boa nova} em meio a ruídos e algoritmos.\\
\textbf{Identidade}: não buscamos seguidores; formamos \textbf{discípulos}.\\
\textbf{Controle}: usamos plataformas como \textbf{meios}, sem entregar o coração a elas.
\end{nicbox}

\section{Estratégia: do público ao pessoal (funil discipulador)}
\begin{enumerate}[leftmargin=*,nosep]
  \item \textbf{Alcance público} (conteúdo aberto) $\rightarrow$ 
  \item \textbf{Conexão pessoal} (DM/e-mail/WhatsApp) $\rightarrow$ 
  \item \textbf{Comunidade} (grupo pequeno/estudo) $\rightarrow$ 
  \item \textbf{Igreja local} (culto, sacramentos, serviço).
\end{enumerate}
Cada etapa tem \textbf{métricas diferentes}: alcance \emph{não} é conversão; conversão \emph{não} é discipulado.

\section{Playbook de evangelismo digital (12 passos que funcionam)}
\begin{checklist}[title={Produção e presença},breakable]
\begin{itemize}[leftmargin=*,noitemsep]
  \item \textbf{Calendário simples}: 2 posts/semana (Palavra + história) e 1 vídeo curto.
  \item \textbf{Conteúdo em 3 formatos}: \emph{Texto} (devocional curto), \emph{Áudio} (1--3 min), \emph{Vídeo} (≤60s, mensagem única).
  \item \textbf{Roteiro 3T} (30--60s): \textbf{Texto} bíblico $\rightarrow$ \textbf{Tradução} à vida $\rightarrow$ \textbf{Tarefa} prática para hoje.
  \item \textbf{Série recorrente}: “Perguntas Difíceis” (cap.~19) em 60s; “Salmos na Prática”; “Histórias de Esperança”.
  \item \textbf{CTA de relacionamento}: sempre convide para DM ou link de inscrição: “podemos orar por você?”
  \item \textbf{Página de aterrissagem} simples: quem somos, horários, grupo no WhatsApp/Telegram, pedido de oração.
  \item \textbf{Moderadores voluntários}: 2--3 pessoas treinadas para acolher DMs, orar e encaminhar.
  \item \textbf{Resposta em 24h}: DM não respondida \emph{esfria}. Defina turnos.
  \item \textbf{Rotina de live} (mensal): Q\&A ao vivo com tema claro e hora fixa.
  \item \textbf{Parcerias locais}: grave testemunhos e projetos de serviço no bairro; mostre \emph{o bem} acontecendo.
  \item \textbf{Próximo passo claro}: estudo bíblico on-line \emph{curto} (4 encontros) que leva ao grupo presencial.
  \item \textbf{Desligamentos saudáveis}: não dependa de uma plataforma; duplique em e-mail/WhatsApp e site simples.
\end{itemize}
\end{checklist}

\section{Roteiros de conteúdo (prontos para gravar)}
\subsection*{A) Devocional de 60 segundos (3T)}
\begin{enumerate}[leftmargin=*,nosep]
  \item \textbf{Texto}: ``Lâmpada para os meus pés é a tua palavra'' (Sl~119:105).
  \item \textbf{Tradução}: quando a luz cai (crise, ansiedade), \emph{primeiro} abrimos a Bíblia.
  \item \textbf{Tarefa}: hoje, leia João 14 e sublinhe um verbo de Jesus; comente aqui o que te deu paz.
\end{enumerate}
\subsection*{B) Testemunho em 3 cenas}
Passado (30s) $\rightarrow$ Encontro com Cristo (30s) $\rightarrow$ Presente (30s com convite: “posso orar por você?”).

\section{Segurança de dados e proteção de pessoas (não pule esta parte)}
\begin{checklist}[title={Padrões mínimos para times de evangelismo}]
\begin{itemize}[leftmargin=*,noitemsep]
  \item \textbf{Consentimento}: peça permissão explícita para armazenar pedidos de oração/dados; permita exclusão fácil.
  \item \textbf{Dados sensíveis}: \emph{não} compartilhe pedidos de oração com detalhes identificáveis em grupos abertos.
  \item \textbf{Contas e acessos}: 2FA em todas as contas; senhas fora de folhas de papel/prints; gestor de senhas.
  \item \textbf{Papéis e logs}: quem responde DMs; quem tem acesso à planilha; \emph{registre} atendimentos sensíveis.
  \item \textbf{Risco de perseguição}: proteja identidades de novos convertidos em contextos hostis; use grupos fechados/anonimizados.
  \item \textbf{Crianças/adolescentes}: regras especiais (sem DMs 1–1; use conta dos pais; conteúdo próprio para faixa etária).
\end{itemize}
\end{checklist}

\section{Discipulado sob pressão (on-line $\rightarrow$ off-line)}
\begin{checklist}[title={Caminho em 4 encontros},breakable]
\begin{itemize}[leftmargin=*,noitemsep]
  \item \textbf{1. Quem é Jesus?} Evangelho claro; convite a ler Marcos em 7 dias.
  \item \textbf{2. Como lemos a Bíblia?} NIC aplicado: \textbf{Narrativa} de Deus, \textbf{Identidade} em Cristo, \textbf{Controle} do pecado.
  \item \textbf{3. Vida nova em comunidade} Batismo, ceia, igreja local, serviço.
  \item \textbf{4. Próximo passo} Ingresso em grupo presencial; acompanhamento 1--1 por 3 meses.
\end{itemize}
\end{checklist}

\section{Métricas que importam (e as que não importam)}
\begin{itemize}[leftmargin=*,noitemsep]
  \item \textbf{Importam}: respostas a DMs, encontros marcados, presença em grupos, decisões por Cristo, batismos, integração em ministérios.
  \item \textbf{Não importam tanto}: views isoladas, curtidas sem conversa, seguidores sem contexto.
\end{itemize}

\section{Equipe mínima (roles e revezos)}
\begin{checklist}
\begin{itemize}[leftmargin=*,noitemsep]
  \item \textbf{Coordenação}: pauta semanal, calendário, métricas simples.
  \item \textbf{Conteúdo}: grava e edita curtos (celular + microfone de lapela).
  \item \textbf{Acolhimento}: responde DMs com empatia e oração.
  \item \textbf{Proteção de dados}: cuida de acessos, consentimentos e planilhas.
\end{itemize}
\end{checklist}

\section{Quando vier a pressão (plataformas limitadas, censura, crises)}
\begin{checklist}[title={Planos B/C}]
\begin{itemize}[leftmargin=*,noitemsep]
  \item \textbf{B}: migre conversas para e-mail/WhatsApp com consentimento; publique em site próprio simples.
  \item \textbf{C}: grupos pequenos presenciais; materiais impressos; rádio local; visitas de casa em casa.
\end{itemize}
\end{checklist}

\section{Perguntas para grupos pequenos}
\begin{itemize}[leftmargin=*,noitemsep]
  \item Quem você pode alcançar \textbf{esta semana} com um vídeo de 60s?
  \item Que \textbf{métrica de discipulado} sua igreja começará a acompanhar mensalmente?
  \item Qual o seu \textbf{plano B} se sua conta ou plataforma cair amanhã?
\end{itemize}

\bigskip
\noindent\textit{Próximo capítulo: Juízos, Babilônia e a queda dos ídolos — o que é não negociável na leitura de Apocalipse.}
